\documentclass[11pt,a4paper]{article}

\usepackage[a4paper,margin=1in]{geometry}
\usepackage{amsmath}
\usepackage{amssymb}
\usepackage{amsthm}
\usepackage{mathpartir}
\usepackage{stmaryrd}
\usepackage{xcolor}
\usepackage{hyperref}
\usepackage{tikz}
\usepackage{tikz-cd}

% Define CCHM-style macros (simplified, avoiding conflicts)
\newcommand{\II}{\mathbb{I}}
\newcommand{\UU}{\mathcal{U}}
\newcommand{\Path}[3]{\mathsf{Path}_{#1}\,#2\,#3}
\newcommand{\comp}[5]{\mathsf{comp}^{#1}\,#2\,[#3 \mapsto #4]\,#5}
\newcommand{\transp}[3]{\mathsf{transp}^{#1}\,#2\,#3}
\newcommand{\fillop}[5]{\mathsf{fill}^{#1}\,#2\,[#3 \mapsto #4]\,#5}
\newcommand{\Glue}[2]{\mathsf{Glue}[#1]\,#2}
\newcommand{\Equiv}[2]{\mathsf{Equiv}\,#1\,#2}
\newcommand{\emptyf}{0_F}
\newcommand{\id}{\mathsf{id}}

\title{Cubical Type Theory Visualizer\\Design Document}
\author{Oguri Cap}
\date{\today}

\begin{document}

\maketitle

\section{Goal}

Transform the current geometric primitive visualizations into \textbf{interactive rule-based visualizations} that teach the core typing rules and computational behavior of Cubical Type Theory.

\section{Current State}

\begin{itemize}
\item[\checkmark] Basic geometric primitives: interval, paths, squares, cubes
\item[$\times$] No typing rules shown
\item[$\times$] No computational behavior demonstrated
\item[$\times$] Limited educational value
\end{itemize}

\section{Design Philosophy}

\subsection{What Makes a Good Rule Visualization?}

\begin{enumerate}
\item \textbf{Show the typing judgment} -- display the rule in CCHM notation
\item \textbf{Geometric intuition} -- 3D representation of what the rule means spatially
\item \textbf{Interactive transformation} -- let users manipulate terms and see how they evaluate
\item \textbf{Step-by-step computation} -- show reduction/evaluation steps
\end{enumerate}

\subsection{Notation System}

Follow CCHM paper exactly:
\begin{itemize}
\item Interval: $\II$ with endpoints $0$, $1$, names $i, j, k$
\item Face formulas: $\varphi, \psi ::= 0_F \mid 1_F \mid (i = 0) \mid (i = 1) \mid \varphi \wedge \psi \mid \varphi \vee \psi$
\item Path types: $\Path{A}{u}{v}$ (dependent paths from $u$ to $v$ in type $A$)
\item Composition: $\comp{i}{A}{\varphi}{u}{a_0}$
\item Transport: $\transp{i}{A}{a}$ (special case: $\comp{i}{A}{\emptyf}{}{a}$)
\item Glue: $\Glue[\varphi \mapsto (T, f)]{A}$
\end{itemize}

\section{Priority Rules to Visualize}

\subsection{Tier 1: Foundation (Must Have)}

These unlock understanding of the entire system.

\subsubsection{Path Types -- Introduction \& Elimination}

\textbf{Typing Rule:}
\begin{mathpar}
\inferrule
  {\Gamma, i : \II \vdash t : A}
  {\Gamma \vdash \langle i \rangle t : \Path{A}{t(i_0)}{t(i_1)}}
\end{mathpar}

\textbf{Visualization:}
\begin{itemize}
\item Show interval $\II$ as a line segment with endpoints
\item Animate $t$ varying along $i$ from $0$ to $1$
\item Display the path as a continuous curve in the type space
\item Interactive: user can scrub $i$ to see $t(i)$ at different points
\end{itemize}

\textbf{Computation:}
\begin{itemize}
\item Path application: $(\langle i \rangle t)\ r$ reduces to $t(i/r)$
\item Show substitution visually
\end{itemize}

\subsubsection{Composition -- The Heart of Cubical}

\textbf{Typing Rule:}
\begin{mathpar}
\inferrule
  {\Gamma, i : \II \vdash A \\ \Gamma \vdash \varphi : F \\ \Gamma, \varphi, i : \II \vdash u : A \\ \Gamma \vdash a_0 : A(i_0)[\varphi \mapsto u(i_0)]}
  {\Gamma \vdash \comp{i}{A}{\varphi}{u}{a_0} : A(i_1)[\varphi \mapsto u(i_1)]}
\end{mathpar}

\textbf{Visualization:}
\begin{itemize}
\item Show an ``open box'' (partial cube with some faces specified by $\varphi$)
\item The base $a_0$ at $i=0$
\item The partial sides $u$ defined on extent $\varphi$
\item Animate filling the box to compute the lid at $i=1$
\item \textbf{Key insight:} ``Being extensible is preserved along paths''
\end{itemize}

\textbf{Example: Transitivity}

Given $p : \Path{A}{a}{b}$ and $q : \Path{A}{b}{c}$, we can construct:
\[
\langle i \rangle \comp{j}{A}{(i=0) \mapsto a, (i=1) \mapsto q\ j}{}{p\ i} : \Path{A}{a}{c}
\]

Visualize as a square with dashed diagonal:
\begin{center}
\begin{tikzcd}
a \arrow[r] \arrow[d, "p"'] & c \\
b \arrow[ur, dashed, "{\text{result}}"'] & 
\end{tikzcd}
\end{center}

\subsubsection{Transport}

\textbf{Typing Rule:}
\begin{mathpar}
\inferrule
  {\Gamma, i : \II \vdash A \\ \Gamma \vdash a : A(i_0)}
  {\Gamma \vdash \transp{i}{A}{a} : A(i_1)}
\end{mathpar}

\textbf{Visualization:}
\begin{itemize}
\item Show type $A$ varying along $i$ (e.g., morphing shape)
\item Element $a$ at $A(i_0)$ being ``transported'' to $A(i_1)$
\item This is $\comp{i}{A}{\emptyf}{}{a}$ (composition with empty extent)
\end{itemize}

\textbf{Example:} Transport along $\Path{\UU}{\mathbb{N}}{\mathbb{B}}$ (if we had a universe)

\subsection{Tier 2: Advanced (Should Have)}

\subsubsection{Kan Filling}

\textbf{Definition:}
\[
\fill{i}{A}{\varphi}{u}{a_0} = \comp{j}{A(i/i \wedge j)}{\varphi \mapsto u(i/i \wedge j), (i=0) \mapsto a_0}{}{a_0}
\]

\textbf{Visualization:}
\begin{itemize}
\item Show both the lid (composition) AND the interior (filling)
\item Animate the filling process from $i=0$ to $i=1$
\item Display: $v(i_0) = a_0$, $v(i_1) = \comp{i}{A}{\varphi}{u}{a_0}$, $v = u$ on $\varphi$
\end{itemize}

\subsubsection{Glue Types -- Univalence}

\textbf{Typing Rule:}
\begin{mathpar}
\inferrule
  {\Gamma \vdash \varphi : F \\ \Gamma, \varphi \vdash T \\ \Gamma, \varphi \vdash f : \Equiv{T}{A} \\ \Gamma \vdash A}
  {\Gamma \vdash \Glue[\varphi \mapsto (T, f)]{A}}
\end{mathpar}

\textbf{Visualization:}
\begin{itemize}
\item Show partial type $T$ on extent $\varphi$
\item Show total type $A$
\item Equivalence $f : T \simeq A$ as a bidirectional morphing
\item Glue operation ``glues'' them together along $\varphi$
\end{itemize}

\textbf{Diagram (from paper):}
\begin{center}
\begin{tikzcd}
& \Glue[\varphi \mapsto (T, f)]{A} \arrow[dl] \arrow[dr] & \\
T \arrow[rr, "f \simeq"] & & A
\end{tikzcd}
\end{center}

\textbf{Example: Path from Equivalence}

Given $f : \Equiv{A}{B}$, define:
\[
E = \Glue[(i=0) \mapsto (A, f), (i=1) \mapsto (B, \id_B)]{B}
\]
Then:
\begin{itemize}
\item At $i=0$: $E = A$
\item At $i=1$: $E = B$
\item Result: $\Path{\UU}{A}{B}$ (univalence!)
\end{itemize}

\subsection{Tier 3: Nice to Have}

\subsubsection{Composition for Specific Types}

Show how composition is defined by induction on types:

\paragraph{Product types:} $C = (x : A) \to B$

For $\Gamma, \varphi, i : \II \vdash \mu : C$ and $\Gamma \vdash \lambda_0 : C(i_0)[\varphi \mapsto \mu(i_0)]$, given $\Gamma \vdash u_1 : A(i_1)$:
\begin{align*}
w &= \fill{i}{A(i/1-i)}{\emptyf}{}{u_1} \\
v &= w(i/1-i) \\
(\comp{i}{C}{\varphi}{\mu}{\lambda_0})\ u_1 &= \comp{i}{B(x/v)}{\varphi \mapsto \mu\ v}{}{\lambda_0\ v(i_0)}
\end{align*}

\paragraph{Path types:} $C = \Path{A}{u}{v}$

For $\Gamma, \varphi, i : \II \vdash p : C$ and $\Gamma \vdash p_0 : C(i_0)[\varphi \mapsto p(i_0)]$:
\[
\comp{i}{C}{\varphi}{p}{p_0} = \langle j \rangle \comp{i}{A}{\varphi \mapsto p\ j, (j=0) \mapsto u, (j=1) \mapsto v}{}{p_0\ j}
\]

\subsubsection{Systems and Face Formulas}

Interactive tool to:
\begin{itemize}
\item Build face formulas $\varphi$ visually (select faces of a cube)
\item Show systems $[\varphi_1 \mapsto t_1, \ldots, \varphi_n \mapsto t_n]$ as partial definitions
\item Visualize compatibility conditions: $\Gamma, \varphi_i \wedge \varphi_j \vdash t_i = t_j$
\end{itemize}

\section{Implementation Strategy}

\subsection{Phase 1: Refactor Current Code}
\begin{enumerate}
\item Extract common 3D utilities (camera, controls, text rendering)
\item Create a \texttt{Rule} component that displays:
  \begin{itemize}
  \item Rule name
  \item Typing judgment (rendered LaTeX or formatted text)
  \item 3D scene
  \item Interactive controls
  \end{itemize}
\item Separate geometric primitives into reusable modules
\end{enumerate}

\subsection{Phase 2: Implement Tier 1 Rules}
\begin{enumerate}
\item \textbf{Path types} -- simplest, builds on existing interval viz
\item \textbf{Transport} -- shows type-varying behavior
\item \textbf{Composition} -- most complex, most important
\end{enumerate}

\subsection{Phase 3: Add Interactivity}
\begin{itemize}
\item Sliders for interval variables ($i$, $j$)
\item Buttons to step through computation
\item Input fields to change terms
\item ``Play'' animation for automatic stepping
\end{itemize}

\subsection{Phase 4: Implement Tier 2 \& 3}
\begin{itemize}
\item Kan filling
\item Glue types
\item Type-specific compositions
\end{itemize}

\section{Technical Considerations}

\subsection{3D Representation Challenges}
\begin{itemize}
\item \textbf{Abstract types:} How to visualize $A : \Type$? Use abstract shapes (spheres, tori)?
\item \textbf{Higher dimensions:} Cubes are 3D, but we have arbitrary dimensions. Use projections?
\item \textbf{Equivalences:} How to show $f : \Equiv{T}{A}$? Morphing animation?
\end{itemize}

\subsection{UI/UX}
\begin{itemize}
\item Split screen: rule on left, visualization on right
\item Step-by-step mode vs. continuous animation
\item Tooltips explaining each part of the rule
\item ``Why does this work?'' explanations
\end{itemize}

\subsection{Performance}
\begin{itemize}
\item Three.js can handle complex scenes
\item May need LOD (level of detail) for higher-dimensional objects
\item Lazy loading for complex rules
\end{itemize}

\section{Success Criteria}

A learner should be able to:
\begin{enumerate}
\item Understand what composition \emph{does} geometrically
\item See how transport is a special case of composition
\item Grasp why Glue types enable univalence
\item Experiment with different terms and see how they compute
\item Build intuition for why the rules are well-typed
\end{enumerate}

\section{References}

\begin{itemize}
\item CCHM Paper: \textit{Cubical Type Theory: a constructive interpretation of the univalence axiom} \\
  \url{https://arxiv.org/abs/1611.02108}
\item Cubical Agda Paper: \textit{Cubical Agda: A Dependently Typed Programming Language with Univalence and Higher Inductive Types} \\
  \url{https://staff.math.su.se/anders.mortberg/papers/cubicalagda2.pdf}
\item Cubical Agda Documentation: \\
  \url{https://agda.readthedocs.io/en/latest/language/cubical.html}
\item 1lab (Cubical Agda formalization): \\
  \url{https://1lab.dev/}
\end{itemize}

\end{document}
